\documentclass[12pt]{article}
\usepackage{graphicx}
\usepackage{amsmath}
\usepackage{geometry}
\geometry{a4paper, margin=1in}
\usepackage{titlesec}
\usepackage{minted}

% Format sections to match the reference document style
\titleformat{\section}
  {\normalfont\Large\bfseries}{\thesection.}{1em}{}
\titleformat{\subsection}
  {\normalfont\large\bfseries}{\thesubsection}{1em}{}

\begin{document}

\begin{center}
    {\Large \textbf{Quantum Simulation Lab}} \\
    \vspace{0.8cm}
    {\Large Experiment III} \\
    \vspace{0.2cm}
    {\Large \textbf{Deutsch Josza Algorithm}} \\
\end{center}

\vspace{0.5cm}
\noindent Name: Hrishikesh M Jayakrishnan \hfill ID: SC26D054
\vspace{1em}
\hrule

\vspace{1em}
\section{Aim}
\begin{enumerate}
    \item To implement and demonstrate the Deutsch Josza algorithm using a quantum circuit.
    \item To verify the correctness of the Deutsch Josza algorithm in determining whether a given function is constant or balanced using the Qiskit framework.
\end{enumerate}
\section{Resources used}
\begin{itemize}
    \item Python
    \item Qiskit 
    \item Jupyter Notebook environment
\end{itemize}

\section{Theoretical Background}
\subsection{Deutsch Josza Algorithm}
The Deutsch Josza algorithm is a quantum algorithm that solves the problem of determining whether a given function \( f: \{0, 1\}^n \rightarrow \{0, 1\} \) is constant (the same output for all inputs) or balanced (outputs 0 for half of the inputs and 1 for the other half). The algorithm achieves this with a single evaluation of the function, demonstrating an exponential speedup over classical algorithms that require multiple evaluations.
\newpage
\subsection{Steps}
Below are the steps involved in implementing the Deutsch-Josza algorithm for a given function \( f: \{0, 1\} \rightarrow \{0, 1\} \):
\begin{enumerate}
    \item Initialize a quantum circuit with \( n + 1 \) qubits, where \( n \) is the number of input bits for the function \( f \). The first \( n \) qubits are initialized to \( |0\rangle \), and the last qubit is initialized to \( |1\rangle \).
    \item Apply a Hadamard gate to all qubits. This creates a superposition of all possible input states for the function.
    \item Apply the oracle \( U_f \) corresponding to the function \( f \). The oracle flips the sign of the amplitude of the state corresponding to the input that maps to 1 under \( f \).
    \item Apply a Hadamard gate to the first \( n \) qubits again. This transforms the state back into a basis where measurement can reveal information about the function.
    \item Measure the first \( n \) qubits. If all measured bits are 0, then the function is constant; otherwise, it is balanced.
\end{enumerate}


\section{Quantum circuit}
The experiment is conducted following these steps in the Qiskit framework:
\begin{enumerate}
    \item Initialize a quantum circuit with \( n + 1 \) qubits and \( n \) classical bits for measurement.
    \item Apply Hadamard gates to all qubits to create a superposition of states.
    \item Apply the oracle \( U_f \) corresponding to the function \( f \) to the qubits. The oracle encodes the function's output into the phase of the quantum state.
    \item Apply Hadamard gates to the first \( n \) qubits again to interfere the amplitudes and extract information about the function.
    \item Measure the first \( n \) qubits.
\end{enumerate}

\begin{figure}[h!]
    \centering
    \includegraphics[width=1.02\textwidth]{../deutsch_josza_algorithm.png}
    \caption{Quantum circuit for the Deutsch-Josza algorithm}
\end{figure}

\begin{figure}[h!]
    \centering
    \includegraphics[width=0.2\textwidth]{../deutsch_josza_constant.png}
    \caption{Quantum circuit acting as an oracle for a constant function}
\end{figure}

\begin{figure}[h!]
    \centering
    \includegraphics[width=1.02\textwidth]{../deutsch_josza_balanced.png}
    \caption{Quantum circuit acting as an oracle for a balanced function}
\end{figure}

\newpage
\section{Measurement Results and Analysis}
The circuit is executed on a quantum simulator for a single shot, and the measurement results are obtained. The results indicate whether the function is constant or balanced based on the measurement of the first qubit.
\begin{figure}[h!]
    \centering
    % Remove the \rule line below and uncomment the \includegraphics line to insert your image
    \includegraphics[width=0.5\textwidth]{../deutsch_josza_alg_measurement.png}
    % \rule{0.5\textwidth}{0.3\textwidth} % This creates a placeholder box
    \caption{Measurement result for a constant function}
\end{figure}
\begin{figure}[h!]
    \centering
    % Remove the \rule line below and uncomment the \includegraphics line to insert your image
    \includegraphics[width=0.5\textwidth]{../deutsch_josza_balanced_measurement.png}
    % \rule{0.5\textwidth}{0.3\textwidth} % This creates a placeholder box
    \caption{Measurement result for a balanced function \( f(x) = x \) for 50\% of the inputs}
\end{figure}
\newpage
The measurement results show that the circuit correctly identifies the nature of the function based on the measurement of the first qubit. If the result is \( |0\rangle \), the function is constant; if the result is \( |1\rangle \), the function is balanced, in a single evaluation of the function, demonstrating the efficiency of the Deutsch-Josza algorithm over classical methods.

\section{Conclusion}
This experiment implements the Deutsch-Josza algorithm using a quantum circuit in the Qiskit framework. The results obtained from the quantum simulator confirm that the algorithm can determine whether a function is constant or balanced with a single evaluation

\newpage
\section{Appendix}
\subsection{Python Code}
\begin{minted}[fontsize=\small]{python}
from qiskit import QuantumCircuit, QuantumRegister, ClassicalRegister, transpile
from qiskit_aer import AerSimulator
from qiskit.visualization import plot_histogram
import numpy as np

def flip_zero_bits(bit_string):
    for qubit, bit in enumerate(reversed(bit_string)):
        if bit == "0":
            oracle.x(qubit)

def create_oracle(n, function_type="random", seed=None):
    if n < 1:
        raise ValueError("n must be at least 1")
    rng = np.random.default_rng(seed)
    oracle = QuantumCircuit(n + 1, name="Oracle")
    if function_type == "random":
        function_type = rng.choice(["constant", "balanced"])
    if function_type not in ("constant", "balanced"):
        raise ValueError("Choose 'constant', 'balanced', or 'random'")
    if rng.integers(0, 2):
        oracle.x(n)
    if function_type == "constant":
        return oracle, function_type
    on_states = rng.choice(
        2**n,
        size=2**(n - 1),
        replace=False
    )
    for state in on_states:
        bit_string = f"{state:0{n}b}"
        flip_zero_bits(bit_string)
        oracle.mcx(list(range(n)), n)
        flip_zero_bits(bit_string)
    return oracle, function_type

n = 4

input_qubits = QuantumRegister(n, name='input')
output_qubit = QuantumRegister(1, name='output')

cr = ClassicalRegister(n, name='c')
qc = QuantumCircuit(input_qubits, output_qubit, cr)
qc.draw(output='mpl')
plt.show()

qc.h(input_qubits)
qc.x(output_qubit)
qc.h(output_qubit)
qc.barrier()
qc.draw(output='mpl')
plt.show()

oracle, actual_function_type = create_oracle(n, function_type="random")
qc.append(oracle.decompose(), list(range(n + 1)))
qc.barrier()
oracle.draw("mpl")
plt.show()

qc.decompose()
qc.draw(output='mpl')
plt.show()

qc.barrier()
qc.h(input_qubits)
qc.measure(input_qubits, cr)
qc.draw("mpl")
plt.show()

print(f"Actual function type: {actual_function_type}")
aer_sim = AerSimulator()
transpiled_qc = transpile(qc, aer_sim)
job = aer_sim.run(transpiled_qc, shots=1)
result = job.result()
counts = result.get_counts(transpiled_qc)
plot_histogram(counts)
plt.show()

\end{minted}
\end{document}
